\par\medskip
\begin{center}\hypertarget{native-163-article116}{}\pdfbookmark[2]{116.}{article-116}116.\end{center}
Bevor man eine Bahn aus geozentrischen Beobachtungen bestimmt, muss man an diesen, sofern höchste Genauigkeit verlangt wird, gewisse Reduktionen wegen Nutation, Präzession, Parallaxe und Aberration vornehmen. Bei einer gröberen Rechnung kann man diese kleinen Korrekturen vernachlässigen.

Beobachtungen von Planeten und Kometen werden gewöhnlich in scheinbaren Rektaszensionen und Deklinationen angegeben, also auf die scheinbare Lage des Äquators bezogen. Da sich diese Lage wegen Nutation und Präzession ändert und daher von Beobachtung zu Beobachtung verschieden ist, empfiehlt es sich, zunächst die veränderliche Ebene durch eine feste zu ersetzen. Dazu kann man entweder den Äquator in seiner mittleren Lage für eine bestimmte Epoche oder die Ekliptik wählen. Meist wird die Ekliptik verwendet, doch auch der Äquator hat eigene Vorteile, die Beachtung verdienen.

Wählt man die Äquatorebene, so sind die Beobachtungen zunächst von der Nutation zu befreien und anschließend unter Berücksichtigung der Präzession auf eine beliebige Epoche zu reduzieren. Dieses Verfahren entspricht genau demjenigen, mit dem aus dem beobachteten Ort eines Fixsterns seine mittlere Position für eine gegebene Epoche abgeleitet wird, und bedarf daher hier keiner Erläuterung. Wählt man dagegen die Ekliptikebene, so stehen zwei Verfahren zur Verfügung. Aus den um Nutation und Präzession korrigierten Rektaszensionen und Deklinationen lassen sich mithilfe der mittleren Schiefe die Längen und Breiten ableiten, wobei die erhaltenen Längen bereits auf das mittlere Äquinoktium bezogen sind. Bequemer ist es, aus den scheinbaren Rektaszensionen und Deklinationen mithilfe der scheinbaren Schiefe die Längen und Breiten zu berechnen und dann die Längen von Nutation und Präzession zu befreien.

Die den einzelnen Beobachtungen entsprechenden Orte der Erde werden aus den Sonnentafeln berechnet. Sie sind offenbar auf dieselbe Ebene zu beziehen wie die Beobachtungen des Himmelskörpers. Bei der Berechnung der Sonnenlänge wird daher die Nutation vernachlässigt. Anschließend wird diese Länge unter Berücksichtigung der Präzession auf die \SourceLeaf{164}feste Epoche reduziert und um 180 Grad vergrößert. Scheint es der Mühe wert, die Sonnenbreite zu berücksichtigen, so gibt man ihr das entgegengesetzte Vorzeichen. So erhält man den heliozentrischen Ort der Erde. Ist der Äquator als Grundebene gewählt, so lässt sich dieser Ort mithilfe der mittleren Schiefe in Rektaszension und Deklination umrechnen.
